\section{Discussion and Conclusions}
\label{sec:conclusion}

\subsection{What the Results Support, and What They Do Not}
\label{sec:discussion}

\paragraph{Feasibility.} Every flow FFE produced passes the schema check, so ``feasible'' means feasible with respect to the 12-feature schema. We did not replay packets and re-extract features with a flow exporter, and a deployed detector may use features outside the twelve (inter-arrival statistics, flag counts) that the four actions also change.

\paragraph{Comparison with the baselines.} FFE has the highest Real ASR on every corpus, significantly so against all seven baselines (Table~\ref{tab:master_evasion_benchmark}), and it stays ahead of CAPGD in 14 of 15 retrained runs with one tie. The ranking should not be read as dominance. FFE is white-box and the black-box PSO and MOEVA are not; against them it needs 2.4--20 times fewer queries, but that is not a like-for-like cost. Surrogate-based attacks need far fewer target queries than FFE (CAPGD 3.2--16.2 against 14.6--93.7) and come within ten points of it on the two easier corpora (88.5\% and 89.5\%), although they fall to 21.1\% on UNSW-NB15, where the surrogate agrees least with the target (87.6\%) and where FFE's margin is largest. The operational-cost advantage (median 0.0006--5.6 units against 15--53) is partly by construction, and in the normalized $L_1$ distance FFE is not smaller than PSO on two corpora. Finally, boundary snapping, which gives the method its second component, adds at most three points of Real ASR at up to 1.8 times the queries (Section~\ref{sec:ablation}); the success of FFE comes mainly from steering along the Lie generators.

\paragraph{Defenses.} Lie-guided adversarial training halves FFE's success at its training budget but protects far less against a larger budget or another attack, and no model we trained brings evasion below 46\% against every attack (Section~\ref{sec:defense}); the defenses were evaluated on one corpus and one seed, and the masked-feature defense implements the infinite-penalty limit, not a penalized split.

\paragraph{Limits of the method.} FFE needs the split thresholds (Tier~1). Under Tier-2 access it depends on transfer, and flows crafted against one architecture evade the others in 24\%--79\% of the cases on CSE-CIC-IDS2018 (Supplementary Section~S5), with CatBoost the hardest target. The packet realization has not been replayed at line rate; control is restricted to forward traffic, deep packet inspection is out of scope, and rate dilution lowers the intensity of volumetric floods.

\paragraph{Dependence on the target.} The CSE-CIC-IDS2018 XGBoost model splits first on destination port, an immutable feature, in 51 of 60 trees, yet is evaded in 99.5\% of the attacked flows. On CIC-IDS2017 XGBoost is evaded in 99.2\% and LightGBM in 35.1\%, and we cannot say whether LightGBM is more robust or the attack fits its splits badly. UNSW-NB15 is the hardest corpus, although its lossy 12-feature mapping also weakens the detector (88.5\% accuracy).

\paragraph{Threats to validity.} The main comparison uses one training seed; the five-seed repetition covers three attacks. Cohorts contain only flows that the models detect when clean (77.6\% of the CSE-CIC-IDS2018 test attacks), so the figures are conditional on detection. MOEVA and PSO ran on 200 flows and the ablations and defenses on 283 or fewer. CAPGD and the NIDSGAN-style generator are our re-implementations with one small surrogate, the baselines use default budgets, and the white-box tree attacks were not run, so the novelty of FFE against them is untested. The feasibility check has a 5\% tolerance on the coupled rates for all methods. ``Full scale'' means the distributed processed files. Exact duplicates are 20\%--43\% of the valid records and are collapsed, but the random split lets near-duplicates fall on both sides, so accuracies are optimistic, and the 12-feature schema makes them incomparable with all-feature studies.

Each generator has a protocol-compliant packet realization: sender pacing for temporal dilation, padding in protocol fields or ignored application frames (\texttt{SSH\_MSG\_IGNORE}) \cite{rfc4253}, payload-free pure ACK or keep-alive chaff \cite{rfc9293,rfc7323} or zero-length UDP datagrams \cite{rfc768}, and handshake window morphing, with sequence numbers and checksums updated via proxy \cite{rfc791,rfc9293,rfc768}.

\subsection{Concluding Remarks and Future Directions}
\label{sec:concluding_remarks}

This article has measured, at full scale, how vulnerable tree-ensemble NIDS with high benchmark accuracy are to evasion that respects a flow schema. Flow-Field Evasion models four packet-level actions as infinitesimal generators on the flow manifold; trajectories in the continuous Lie cone preserve the conservation invariants and a projection enforces them on the discrete schema. On the complete, deduplicated CSE-CIC-IDS2018, CIC-IDS2017 and UNSW-NB15 captures, FFE evades XGBoost in 99.5\%, 99.2\% and 57.6\% of the detected attacks it targets, with schema-feasible flows and 14.6--93.7 target queries per flow. It is significantly more successful than all seven baselines, including surrogate-based CAPGD, and stays ahead over five retrained seeds, while CAPGD needs far fewer target queries and boundary snapping contributes at most three points. Evasion depends on the target: LightGBM on CIC-IDS2017 resists FFE far better than XGBoost, and CatBoost resists transfer from the other ensembles. Lie-guided adversarial training halves FFE success at the budget it was trained on but protects much less against stronger attacks, and no defense we tested reduces evasion below 46\%. Future work includes a packet-level replay with feature re-extraction, white-box tree-attack controls, defenses that withstand attacks other than the one used for augmentation and are evaluated on all corpora and several seeds, the remaining public NIDS corpora, line-rate replay on a hardware testbed, and certified robustness bounds for adversarially trained ensembles.

